\originalchapter{不可分辨粒子与多体态}\label{ch:states}
\originalsection{交换对称性}
先明确单粒子态怎样组成多体空间。非相互作用时 $H=\sum_i h_i$，若 $h\ket a=\epsilon_a\ket a$，乘积态 $\ket{a_1}\otimes\cdots\otimes\ket{a_N}$ 为带暂时标签粒子的本征态，能量相加。标签 $a$ 同时包含空间与自旋；加入相互作用后这些乘积仍可作展开基底，但一般不再是总 Hamiltonian 的本征态。同种粒子须再把此大空间限制到正确的交换子空间。
两个同种粒子交换名字不会产生可观测的新构型. 在通常三维量子体系中, 整数自旋粒子采用对称Bose子空间, 半整数自旋粒子采用反对称Fermi子空间. 这是自旋--统计规律的物理输入; 单凭概率密度在交换下不变不足以推导完整定理, 也不能替代相对论场论证明. 单粒子标签必须包括空间态和自旋态, 两个电子可以共享空间轨道而自旋不同.

正交单粒子态$a,b$构成的两体乘积$a(1)b(2)$通常不满足所需交换性. 不同态时, 对称或反对称组合为$[a(1)b(2)\pm b(1)a(2)]/\sqrt2$. 若$a=b$, 反对称组合为零, 给Pauli不相容; 对称态则只是$a(1)a(2)$, 不能继续沿用$1/\sqrt2$的归一化因子.

对$N$个粒子, 将全部置换$P$的乘积态相加, Bose用正号, Fermi用置换奇偶性$(-1)^P$. 若某单粒子态占$n_\alpha$次, Bose全部置换和的归一化分母为$\sqrt{N!\prod_\alpha n_\alpha!}$. 因为有$N!$个排列, 内积不为零的成对排列有$\prod n_\alpha!$种重复. Fermi各$n_\alpha$只能是0或1, 分母退为$\sqrt{N!}$. 这解释了占据数表示的来历, 不是给每个不可分辨粒子分配可观测编号.

\originalsection{占据数与计数差别}
理想多体能量为$E=\sum_\alpha n_\alpha\epsilon_\alpha$, 粒子数$N=\sum_\alpha n_\alpha$. 给定占据数就指定一个对称性允许的多体态. 相同$E$可能有许多不同占据数组合, 仍需分别计入简并度. 对可区分粒子, 将$N$个粒子分配到这些态时还有$N!/\prod n_\alpha!$个排列权重; Bose/Fermi态不再包含这种标签排列.

若两个无自旋粒子在三个等能单粒子态中, Bose多体态数为$\binom{3+2-1}2=6$, Fermi态数$\binom32=3$, 可区分粒子态数$3^2=9$. 前两者分别数可重复组合与不可重复组合, 能量相同也不能把统计当作相同.

课内的两粒子量子统计例用相隔能隙$\epsilon$的单粒子能级说明热分布. 若只有两单粒子态, 两个Boson允许$(2,0),(1,1),(0,2)$, 总配分函数为$1+\e^{-\beta\epsilon}+\e^{-2\beta\epsilon}$; 两个无自旋Fermion只允许$(1,1)$, 配分函数为$\e^{-\beta\epsilon}$. 两个可区分粒子中间占据有两种, 配分函数为$(1+\e^{-\beta\epsilon})^2$. 若有自旋简并, 必须先展开为完整单粒子态再计数.

\originalsection{氦与核自旋的实例}
${}^4$He原子有偶数个Fermion组成, 基态总自旋为0, 是Boson; ${}^3$He有半整数总自旋, 是Fermion. 电子间交换也限制氦原子的单重态和三重态: 两电子自旋单重态反对称, 故空间部分对称; 自旋三重态对称, 故空间部分反对称. 基态空间轨道两电子相同, 只能配单重态. 这与两核转动Raman交替强度是同一总波函数对称原则的两次应用.

固定$N$时, 正则和$\sum_{\{n_\alpha\}:\sum n_\alpha=N}\prod_\alpha\e^{-\beta\epsilon_\alpha n_\alpha}$带有总数约束, 一般不能化为每个模式独立求和的乘积. 将固定$N$的计算硬拆为独立模式, 会错计多体态. 下一章通过允许粒子交换移除这个约束, 再用化学势使\textbf{平均}总数达到给定值.
